% Member 4 - final report, IV.H Discussion and Limitations
\subsection{Discussion and Limitations}
\label{sec:limitations}

\textit{What the experiments establish.} Style augmentation is effective on PACS under a single-source protocol: adding 1{,}349 styled images to the 1{,}349 original photographs raises the cross-domain average by several points, and it does so while reducing the seed-to-seed spread. The controls establish that the effect is not a data-volume artefact --- doubling the training rows with byte-identical copies does not reproduce it, and neither does doubling them with non-stylistic transforms. What remains is the style itself.

\textit{What the experiments do not establish.} The spatial modulation does not contribute. Under a strength-matched control the two methods are statistically indistinguishable, and the equal-strength variant confirms this from the opposite direction. The mechanism that motivated the method --- protecting the content's structure where it matters --- therefore does not translate into a classification benefit under this protocol. Two observations explain why this is plausible rather than surprising. First, the structural response map is a soft Sobel magnitude: it responds to any sharp intensity change, including background texture and object edges alike, so it cannot distinguish the class-relevant contour from an irrelevant one. Second, the analysis in Section~\ref{sec:mechanism} shows that the modulation does shift the amount of style that reaches each region, but the resulting per-image edge retention does not improve on the uniform control, so the "protection" is not selective either.

\textit{Limitations.} We state the following limits of the evidence:

\begin{itemize}
\item \textit{One architecture, one dataset.} Every experiment uses ResNet-18 and PACS. The conclusion that the spatial modulation does not help is a statement about this setting; a different architecture, a different dataset, or a stronger stylizer could behave differently. The MixStyle comparison in Section~\ref{sec:ablation} gives one reference point, not a broad survey.
\item \textit{Five seeds.} The main table uses five seeds. That is enough to detect the headline effect and to rule out the larger effects, but the smaller differences --- for instance the spatial modulation's effect --- remain inside the confidence interval, so we can say that they are small, not that they are exactly zero.
\item \textit{A single style pool for the main results.} The main experiments use one frozen pool of ten external paintings. The sensitivity experiment in Section~\ref{sec:ablation} repeats the pipeline with two disjoint halves of that pool, which tests the dependence on the pool's size but not on the choice of a completely different set of paintings.
\item \textit{The decomposition's third row is an upper bound.} As noted in Table~\ref{tab:decomposition}, the contribution attributed to style simultaneously contains the comparison of styling against a generic transform and the comparison of one particular family of transforms against a random one. A finer set of controls would separate these.
\item \textit{The style reference is decoupled from the object.} Every training photo is paired with a painting that has nothing to do with the depicted object. This is what makes the style pool usable without leaking the target domains, but it also means that the generated image can be semantically odd. A method that preserved the object's identity through the stylization might behave differently.
\end{itemize}

\textit{Implications for practice.} The practical recommendation that follows is narrower than "style augmentation works". It is that a style-augmented training set should be compared against a control of the same shape, and that an adaptive variant should be compared against a strength-matched uniform one. Both comparisons are cheap --- the controls are one extra generator and one extra training run per configuration --- and both change the conclusion of the experiment.
